\documentclass[UTF8,a4paper,11pt]{ctexart}

\usepackage[a4paper,margin=2.2cm]{geometry}
\usepackage{amsmath,amssymb}
\usepackage{booktabs}
\usepackage{array}
\usepackage{longtable}
\usepackage{xcolor}
\usepackage{enumitem}
\usepackage{hyperref}
\usepackage{siunitx}
\usepackage{float}
\usepackage{caption}
\usepackage{fancyhdr}
\usepackage{lastpage}

\hypersetup{
  colorlinks=true,
  linkcolor=blue!50!black,
  urlcolor=blue!50!black
}

\captionsetup{font=small,labelfont=bf}
\setlist[itemize]{leftmargin=1.8em}
\setlist[enumerate]{leftmargin=2.0em}
\setlength{\parindent}{2em}
\setlength{\parskip}{0.35em}
\setlength{\headheight}{14pt}

\pagestyle{fancy}
\fancyhf{}
\lhead{qDIC + WLI 标准实验报告模板}
\rhead{\today}
\cfoot{\thepage/\pageref{LastPage}}

\newcommand{\Code}[1]{\texttt{\detokenize{#1}}}
\newcommand{\TODO}[1]{\textcolor{red!70!black}{[TODO: #1]}}

% ---------- Metadata ----------
\newcommand{\ExperimentTitle}{qDIC + 白光干涉仪微透镜联合采集实验}
\newcommand{\ExperimentDate}{\today}
\newcommand{\ExperimentLocation}{实验地点待填写}
\newcommand{\Experimenter}{实验者待填写}
\newcommand{\Partner}{同组人员待填写}
\newcommand{\InstrumentA}{qDIC / 定量微分相衬显微镜}
\newcommand{\InstrumentB}{WLI / 白光干涉仪}
\newcommand{\SampleName}{微透镜样品待填写}
\newcommand{\ReportID}{LAB-YYYYMMDD-001}

\begin{document}

\begin{center}
  {\LARGE \bfseries 标准实验报告}\\[1.0em]
  {\Large \bfseries \ExperimentTitle}\\[1.5em]
\end{center}

\begin{longtable}{>{\bfseries}p{0.22\textwidth} p{0.70\textwidth}}
\toprule
项目 & 内容 \\
\midrule
\endfirsthead
\toprule
项目 & 内容 \\
\midrule
\endhead
实验编号 & \ReportID \\
实验日期 & \ExperimentDate \\
实验地点 & \ExperimentLocation \\
实验者 & \Experimenter \\
同组人员 & \Partner \\
主要仪器 & \InstrumentA；\InstrumentB \\
实验样品 & \SampleName \\
\bottomrule
\end{longtable}

\tableofcontents
\newpage

\section{实验标题}

\ExperimentTitle

\section{实验目的}

\begin{enumerate}
  \item 通过 qDIC 采集微透镜在 \Code{clean / defect} 状态下的真实观测数据。
  \item 通过 WLI 获取同一颗微透镜的三维形貌，用作几何参考。
  \item 建立同一颗微透镜在两种状态下的 qDIC 与 WLI 跨模态配对关系。
  \item 验证真实样本是否满足当前方法所依赖的“平滑标准面 + 稀疏小缺陷”假设。
  \item 为后续真实数据验证、QC 门控以及算法适配提供基础实验数据。
\end{enumerate}

\section{实验原理}

\subsection{qDIC 原理}

qDIC 通过相位梯度或光程差引起的强度变化形成对比，可用于观察微透镜表面局部结构、边缘变化和缺陷引起的微小扰动。
在本实验中，qDIC 作为主观测模态，对应后续算法处理的真实输入。

\subsection{WLI 原理}

白光干涉仪通过干涉扫描获取样品表面的三维高度信息，可给出计量级表面形貌。
在本实验中，WLI 作为几何参考模态，用于：
\begin{itemize}
  \item 建立 clean 状态下的标准面
  \item 提取 defect 状态相对 clean 状态的残差形貌
  \item 判断真实样本是否符合稀疏缺陷假设
\end{itemize}

\subsection{联合实验原理}

对同一颗微透镜，分别在 qDIC 和 WLI 下采集 \Code{clean} 与 \Code{defect} 两种状态的数据。
通过跨模态配对，将 qDIC 的真实观测与 WLI 的几何参考对应起来，从而实现：
\begin{itemize}
  \item 真实输入与几何参考的联合分析
  \item clean / defect 状态差异的局部化解释
  \item 对当前主线假设的直接验证
\end{itemize}

\section{实验仪器与材料}

\subsection{实验仪器}

\begin{longtable}{>{\bfseries}p{0.26\textwidth} p{0.22\textwidth} p{0.42\textwidth}}
\toprule
仪器名称 & 型号/规格 & 备注 \\
\midrule
\endfirsthead
\toprule
仪器名称 & 型号/规格 & 备注 \\
\midrule
\endhead
qDIC 显微镜 & \TODO{填写} & 主观测模态 \\
WLI 白光干涉仪 & \TODO{填写} & 三维形貌参考 \\
物镜 & \TODO{填写} & 倍率、NA \\
样品台 / 载台 & \TODO{填写} & 建议可记录 stage 坐标 \\
数据采集工作站 & \TODO{填写} & 记录软件版本与导出格式 \\
\bottomrule
\end{longtable}

\subsection{实验材料}

\begin{itemize}
  \item 微透镜样品：\TODO{填写样品编号、批次、数量}
  \item 若进行缺陷注入：\TODO{填写注入材料或工具，例如微粒、附着物、划痕工具}
  \item 辅助定位材料：\TODO{方向标记、fiducial、样品编号标签等}
\end{itemize}

\section{实验步骤}

\subsection{实验前准备}

\begin{enumerate}
  \item 检查 qDIC 与 WLI 两台仪器状态，完成开机与预热。
  \item 完成 qDIC 的平场、暗场或参考件校准。
  \item 完成 WLI 的当天标定与参考检查。
  \item 确定样品方向标记、样品编号和 overview 成像方案。
  \item 记录当天环境条件、仪器参数、操作者和软件版本。
\end{enumerate}

\subsection{clean 状态采集}

\begin{enumerate}
  \item 固定样品方向，在低倍视场下确定目标微透镜位置。
  \item 切换到目标视场，保证视场中仅包含一颗微透镜且 aperture 完整。
  \item 对 clean 状态样品进行 qDIC 重复采集，记录 repeat 编号与 stage 坐标。
  \item 对同一颗 clean 微透镜进行 WLI 扫描，保存高度图与有效 mask。
\end{enumerate}

\subsection{缺陷注入与 defect 状态采集}

\begin{enumerate}
  \item 对目标样品引入可控缺陷，记录缺陷类型、位置、方式和是否可逆。
  \item 返回同一颗微透镜，重新对齐视场。
  \item 对 defect 状态进行 qDIC 重复采集。
  \item 对同一颗 defect 微透镜进行 WLI 扫描。
\end{enumerate}

\subsection{实验后检查}

\begin{enumerate}
  \item 检查 qDIC 和 WLI 原始文件是否完整且可读。
  \item 检查 clean / defect 标记是否一致。
  \item 检查同一颗 lens 在 qDIC 与 WLI 中是否成功配对。
  \item 检查 clean WLI residual 是否显著小于 defect residual。
\end{enumerate}

\section{数据记录}

\subsection{实验条件记录表}

\begin{longtable}{>{\bfseries}p{0.24\textwidth} p{0.30\textwidth} p{0.30\textwidth}}
\toprule
参数 & qDIC & WLI \\
\midrule
\endfirsthead
\toprule
参数 & qDIC & WLI \\
\midrule
\endhead
仪器型号 & \TODO{填写} & \TODO{填写} \\
物镜倍率 & \TODO{填写} & \TODO{填写} \\
NA & \TODO{填写} & \TODO{填写} \\
波长 / 光源 & \TODO{填写} & \TODO{填写} \\
像素尺寸 / lateral sampling & \TODO{填写} & \TODO{填写} \\
FOV & \TODO{填写} & \TODO{填写} \\
曝光时间 & \TODO{填写} & \TODO{填写} \\
增益 & \TODO{填写} & \TODO{填写} \\
扫描范围 / z range & \TODO{若适用} & \TODO{填写} \\
\bottomrule
\end{longtable}

\subsection{样本采集记录表}

\begin{longtable}{>{\bfseries}p{0.14\textwidth} p{0.10\textwidth} p{0.10\textwidth} p{0.10\textwidth} p{0.14\textwidth} p{0.26\textwidth}}
\toprule
样品编号 & 状态 & 模态 & 重复号 & 配对样本 & 原始文件路径 / 备注 \\
\midrule
\endfirsthead
\toprule
样品编号 & 状态 & 模态 & 重复号 & 配对样本 & 原始文件路径 / 备注 \\
\midrule
\endhead
\TODO{lens\_a} & clean & qDIC & 1 & \TODO{wli\_lens\_a\_clean\_r01} & \TODO{填写} \\
\TODO{lens\_a} & clean & WLI & 1 & \TODO{qdic\_lens\_a\_clean\_r01} & \TODO{填写} \\
\TODO{lens\_a} & defect & qDIC & 1 & \TODO{wli\_lens\_a\_defect\_r01} & \TODO{填写} \\
\TODO{lens\_a} & defect & WLI & 1 & \TODO{qdic\_lens\_a\_defect\_r01} & \TODO{填写} \\
\bottomrule
\end{longtable}

\subsection{元数据文件}

推荐同时保留结构化 manifest：
\begin{itemize}
  \item \Code{/home/void0312/PINNs/mini_grin_rebuild/external_data/templates/qdic_wli_session_manifest.example.json}
\end{itemize}

\section{数据处理}

\subsection{qDIC 数据处理}

\begin{enumerate}
  \item 导入原始 qDIC 数据。
  \item 做必要的平场、暗场、强度归一化与坏点处理。
  \item 选择单颗微透镜视场，裁剪有效 aperture 区域。
  \item 对 clean / defect 配对样本进行差分或重复性分析。
\end{enumerate}

\subsection{WLI 数据处理}

\begin{enumerate}
  \item 导入 WLI 高度图和有效 mask。
  \item 去除 tilt / piston / form。
  \item 对 clean 状态拟合标准面。
  \item 计算 defect 状态相对 clean 标准面的残差形貌。
\end{enumerate}

\subsection{跨模态处理}

\begin{enumerate}
  \item 建立 qDIC 与 WLI 的 lens identity 对应关系。
  \item 依据 stage 坐标、方向标记、overview 图完成跨模态配对。
  \item 统一 clean / defect 状态定义与文件命名。
  \item 输出后续算法可用的配对数据单元。
\end{enumerate}

\section{结果与数据分析}

\subsection{重复性分析}

\TODO{记录 clean 状态下 qDIC 重复采集与 WLI 重复扫描的一致性结果。}

\subsection{标准面与缺陷残差分析}

\TODO{写明 clean WLI 标准面拟合结果，以及 defect 相对 clean 的残差是否局部且稀疏。}

\subsection{跨模态配对结果}

\TODO{写明 qDIC 与 WLI 是否成功完成同一颗 lens 的 clean / defect 配对。}

\subsection{数据处理结果表}

\begin{longtable}{>{\bfseries}p{0.16\textwidth} p{0.18\textwidth} p{0.18\textwidth} p{0.18\textwidth} p{0.20\textwidth}}
\toprule
样品编号 & qDIC 重复性 & WLI 重复性 & 缺陷残差特征 & 是否满足主线假设 \\
\midrule
\endfirsthead
\toprule
样品编号 & qDIC 重复性 & WLI 重复性 & 缺陷残差特征 & 是否满足主线假设 \\
\midrule
\endhead
\TODO{lens\_a} & \TODO{填写} & \TODO{填写} & \TODO{填写} & \TODO{是/否} \\
\TODO{lens\_b} & \TODO{填写} & \TODO{填写} & \TODO{填写} & \TODO{是/否} \\
\bottomrule
\end{longtable}

\section{误差分析}

建议从以下几类误差展开：

\begin{enumerate}
  \item \textbf{仪器误差}：qDIC 光路漂移、WLI 标定误差、焦点或扫描不稳定。
  \item \textbf{定位误差}：两台仪器下未能精确回到同一颗微透镜。
  \item \textbf{样品误差}：样品方向变化、污染扩散、缺陷注入不可控。
  \item \textbf{处理误差}：去 tilt / 去 form 过强或不足，导致残差被放大或掩盖。
  \item \textbf{配对误差}：clean / defect 或 qDIC / WLI 文件命名、编号、状态对应错误。
\end{enumerate}

\TODO{在此补充本次实验实际观察到的主要误差来源。}

\section{问题讨论}

建议围绕以下问题展开：

\begin{enumerate}
  \item qDIC 与 WLI 是否足以构成稳定的同 lens 配对实验？
  \item 真实微透镜是否普遍满足“平滑标准面 + 稀疏小缺陷”假设？
  \item 哪类样本明显超出当前方法适用范围？
  \item 当前实验流程中最薄弱的环节是采集、配准、还是缺陷注入？
  \item 下一轮应优先扩大样本量，还是先优化实验流程？
\end{enumerate}

\section{实验结论}

\TODO{用 3--6 句话总结实验结论。}

\section{附录}

\subsection{附录 A：原始文件与路径}

\begin{itemize}
  \item qDIC 原始文件目录：\TODO{填写}
  \item WLI 原始文件目录：\TODO{填写}
  \item manifest 文件路径：\TODO{填写}
\end{itemize}

\subsection{附录 B：后续建议}

\begin{enumerate}
  \item \TODO{建议 1}
  \item \TODO{建议 2}
  \item \TODO{建议 3}
\end{enumerate}

\end{document}
